Flood damage assessment is a cornerstone of physical climate risk quantification for financial institutions. The Basel Committee on Banking Supervision~\cite{basel2021a, basel2021b} identifies data scarcity  --  particularly the lack of paired hazard-loss observations for damage function calibration  --  as the single most critical gap in the physical risk assessment pipeline. Existing depth-damage functions, such as those maintained by the Joint Research Centre~\cite{huizinga2017global} and the US Army Corps of Engineers~\cite{usace2003}, treat damage as an instantaneous function of water depth. These functions underpin widely used platforms including CLIMADA~\cite{aznar2019climada} and form the basis for regulatory stress tests~\cite{ecb2025stress}. Yet they ignore a physically obvious variable: the duration of inundation. A flood that recedes in hours causes fundamentally different damage from one that persists for weeks, through material saturation, mould growth, electrical system failure, and structural weakening.

This omission has measurable consequences. Using 295{,}637 cleaned FEMA NFIP claims from three major US hurricanes, we show that standard JRC depth-damage functions systematically underpredict damage by about 15 percentage points for prolonged events -- a bias that propagates directly into credit risk estimates for flood-exposed companies.

\subsection{Contribution}

This paper makes three contributions:

\begin{enumerate}
    \item We extend the generalised logistic (Richards) damage function with a \textit{duration-dependent formulation} in which the curve's inflection point, onset speed, and asymmetry depend continuously on hazard duration. We calibrate the seven parameters jointly on 295{,}637 FEMA NFIP flood insurance claims spanning Hurricanes Sandy (36\,h), Harvey (288\,h), and Katrina (120\,h). The key empirical finding is that duration primarily affects damage through inflection point shift ($\lambda_a = 0.49$), while curve shape parameters are not significantly non-zero in the three-event calibration ($\lambda_Q \approx 0$, $\delta_\nu \approx 0$); this conclusion is, however, limited by the fact that the calibration spans only three distinct event-level duration values.

    \item We introduce a \textit{sector adjustment multiplier} $\kappa$ calibrated empirically from 82{,}650 NFIP non-residential claims, grouped by number of stories and insurance coverage. A coverage-gradient extrapolation bridges NFIP-scale commercial buildings to corporate-scale facilities without relying on corporate loss disclosures. Out-of-sample validation on 24 companies across six flood events yields a median predicted-to-actual ratio of $0.90\times$ (23/24 within 0.4--2.5$\times$).

    \item We demonstrate the \textit{full calibrated pipeline} on Archer Daniels Midland's 12 flood-exposed Midwest facilities, producing facility-level loss estimates with duration-dependent damage ratios, CapEx/OpEx decomposition, insurance offsets, and credit risk translation via an interest coverage ratio approach. The framework is modular: the credit risk translation module can employ alternative models (including the Merton structural model~\cite{merton1974pricing}) without changing the damage assessment.
\end{enumerate}

\subsection{Data requirements}

The framework requires: (i) asset locations and replacement values, obtainable from company filings; (ii) local flood hazard parameters (depth, duration), available from FEMA flood maps and hydrological models; (iii) financial statements (EBIT, debt, interest expenses) for credit risk translation. The damage function parameters are provided by our calibration and do not require company-specific fitting.

\subsection{Paper organisation}

Section~\ref{sec: lit review} reviews damage function literature and the calibration data gap. Section~\ref{sec: methods} presents the framework architecture, the duration-dependent Richards formulation, sector adjustment, and credit risk translation. Section~\ref{sec: empirical} describes the FEMA dataset, joint calibration, and four validation tests. Section~\ref{sec: adm} applies the calibrated pipeline to ADM. Section~\ref{sec: event study} presents the Hurricane Harvey event study. Section~\ref{sec: discussion} discusses findings and limitations. Section~\ref{sec: conclusion} concludes.
